The first instruction word contains a four-bit prefix/length area and a twelve-bit primary payload. Compact
instructions place operand fields directly in this primary payload when the field set fits. Extended
instructions use a primary root and a following descriptor/opcode word, with more payload words added only
when needed.

\Needspace{2.15in}
\begin{center}\vspace{3pt}
\begin{tikzpicture}[x=0.300in,y=0.22in,every node/.style={font=\scriptsize}]
\node[anchor=east] at (-0.35,0.50) {primary};
\node[anchor=south] at (0,1.18) {15};
\node[anchor=south] at (16,1.18) {0};
\draw[LightRule] (4,-0.00) -- (4,1.00);
\draw[LightRule] (8,-0.00) -- (8,1.00);
\draw[LightRule] (12,-0.00) -- (12,1.00);
\draw (0,-0.00) rectangle (4,1.00);
\node at (2.00,0.50) {\texttt{----}};
\draw (4,-0.00) rectangle (16,1.00);
\node at (10.00,0.50) {\texttt{p}};
\node[anchor=east] at (-0.35,-1.58) {descriptor};
\node[anchor=south] at (0,-0.90) {15};
\node[anchor=south] at (16,-0.90) {0};
\draw[LightRule] (4,-2.08) -- (4,-1.08);
\draw[LightRule] (8,-2.08) -- (8,-1.08);
\draw[LightRule] (12,-2.08) -- (12,-1.08);
\draw (0,-2.08) rectangle (16,-1.08);
\node at (8.00,-1.58) {\texttt{s}};
\node[anchor=east] at (-0.35,-3.66) {payload};
\node[anchor=south] at (0,-2.98) {15};
\node[anchor=south] at (16,-2.98) {0};
\draw[LightRule] (4,-4.16) -- (4,-3.16);
\draw[LightRule] (8,-4.16) -- (8,-3.16);
\draw[LightRule] (12,-4.16) -- (12,-3.16);
\draw (0,-4.16) rectangle (16,-3.16);
\node at (8.00,-3.66) {\texttt{x}};
\end{tikzpicture}
\manualfigurecaption{Instruction Word Format Families}
\end{center}


@REPRESENTATIVE_ENCODINGS@
